\begin{textsample}{NEG Dim 2 – persona\_gemini – Score 1.00 – t870\_gemini.txt}  \label{ex:f2_neg_008}
A Greenpeace China campaign, which started in 2008 amid food scandals, has led to significant developments in food safety.
Recent tests on produce from 6 \textbf{supermarkets} in 8 cities, however, showed that 87 % of samples had excessive pesticides.
The highest levels were found on plastic-wrapped items.
Our campaign included a 2015 lawsuit against the \textbf{retailer} Lianhua.
This action prompted Shanghai to establish a city-wide traceability system for fresh produce, which has become a national model for improving food safety.
As part of this system, the city sources 50 % of its produce locally.
The campaign also resulted in Lianhua dropping 46 harmful pesticides.
At a national level, China ’s 2015 policy set a goal of achieving zero increase in pesticide use by 2020.
Sustainable organic farming is emerging as an alternative approach.

% matched lemmas: retailer, supermarket
\end{textsample}
